\chapter{Tests, Déploiement et Étude Technico-Économique}
\label{chap:tests-deploiement}

\section{Introduction}
Ce chapitre valide la plateforme \textit{RohWinBghit} à travers une campagne de tests automatisés, une évaluation d'utilisabilité et des mesures de performance. Il présente ensuite la stratégie de déploiement, puis l'étude technico-économique dans le cadre de l'Arrêté~1275.

% ================================================================
%  PARTIE I --- TESTS ET VALIDATION
% ================================================================

\section{Tests de l'application}
\label{sec:tests-app}

La validation repose sur une stratégie combinant tests automatisés Jest et campagnes de tests manuels. Le tableau~\ref{tab:tests} synthétise les résultats du 24~avril~2026.

\begin{table}[H]
\centering
\footnotesize
\begin{tabularx}{\textwidth}{|X|c|c|c|}
\hline
\rowcolor{rohlight}
\textbf{Module} & \textbf{Cas} & \textbf{Réussis} & \textbf{Échoués} \\
\hline
\multicolumn{4}{|l|}{\cellcolor{rohlight!40}\textit{Tests unitaires (logique pure)}} \\
\hline
Authentification, \gls{jwt}, \textit{blacklist} & 50 & 50 & 0 \\
\hline
Tarification, heatmap, score conducteur & 16 & 16 & 0 \\
\hline
Détection de fraude \textsc{kyc} & 6 & 6 & 0 \\
\hline
Machine d'états des trajets \& \textit{matching} & 7 & 7 & 0 \\
\hline
Utilitaires \& \textit{middlewares} & 28 & 28 & 0 \\
\hline
\rowcolor{rohlight!20}
\textbf{Sous-total unitaires} & \textbf{107} & \textbf{107} & \textbf{0} \\
\hline
\multicolumn{4}{|l|}{\cellcolor{rohlight!40}\textit{Tests d'intégration (\textsc{postgresql} + HTTP + Socket.IO)}} \\
\hline
Trajets~: recherche, réservation, gestion & 20 & 18 & 2 \\
\hline
Messagerie temps réel & 14 & 13 & 1 \\
\hline
Utilisateurs \& profils & 6 & 5 & 1 \\
\hline
Modèles de réservation & 4 & 4 & 0 \\
\hline
\rowcolor{rohlight!20}
\textbf{Sous-total intégration} & \textbf{44} & \textbf{40} & \textbf{4} \\
\hline
\rowcolor{rohlight}
\textbf{Total général} & \textbf{151} & \textbf{147} & \textbf{4} \\
\hline
\end{tabularx}
  \caption{Synthèse des tests automatisés Jest (24 avril 2026)}
  \label{tab:tests}
\end{table}

Les 107~tests unitaires passent à 100\,\%. Sur les 44~tests d'intégration, 40 réussissent (90,9\,\%). Les 4~échecs résiduels sont localisés dans les scénarios de synchronisation WebSocket et de matching spatial asynchrone~; ils traduisent des difficultés d'isolation transactionnelle de l'environnement de test, non des défauts fonctionnels. Le taux de réussite global est de \textbf{97,4\,\%} (147/151), satisfaisant l'objectif ENF~5 ($\geq$~80\,\%). La stabilisation du harnais via \textit{Testcontainers} est un chantier prioritaire post-mémoire.

\subsection{Scénarios de tests fonctionnels}
\label{sec:scenarios-tests}

\begin{table}[H]
\centering
\footnotesize
\renewcommand{\arraystretch}{1.3}
\begin{tabularx}{\textwidth}{|l|p{2.5cm}|p{3.5cm}|p{3.5cm}|c|}
\hline
\rowcolor{rohlight}
\textbf{ID} & \textbf{Scénario} & \textbf{Étapes} & \textbf{Résultat Attendu} & \textbf{Statut} \\
\hline
TS-01 & Inscription passager & Saisie téléphone $\to$ OTP $\to$ Validation & Compte créé, JWT émis & \textbf{Succès} \\
\hline
TS-02 & Vérification KYC & CIN recto/verso $\to$ Capture faciale $\to$ Pipeline IA & Similarité faciale $>$ 0,60, statut validé & \textbf{Succès} \\
\hline
TS-03 & Publication de trajet & Sélection wilayas $\to$ Date et prix $\to$ Soumission & Trajet enregistré, statut \texttt{active} & \textbf{Succès} \\
\hline
TS-04 & Réservation de place & Recherche $\to$ Sélection $\to$ Réservation & Places diminuées, notification conducteur & \textbf{Succès} \\
\hline
TS-05 & Paiement électronique & Choix Edahabia $\to$ Carte sandbox $\to$ Validation & Débit simulé, statut \texttt{confirmed} & \textbf{Succès} \\
\hline
TS-06 & Messagerie WebSocket & Ouverture chat $\to$ Envoi $\to$ Vérification réception & Message délivré via Socket.IO, persisté & \textbf{Succès} \\
\hline
\end{tabularx}
\caption{Matrice des scénarios de tests fonctionnels}
\label{tab:test-matrix}
\end{table}

\section{Résultats des tests}
\label{sec:resultats-tests}

\subsection{Évaluation d'utilisabilité --- échelle SUS}

L'échelle SUS~\cite{brooke_sus_1996} agrège dix affirmations en un score $\in [0, 100]$~; un score $\geq$~68 est considéré comme «~bon~»~\cite{Bangor2008}. Une session de bêta-test a été conduite auprès de \textbf{18~participants} (12~étudiants de la Faculté des Sciences de Tlemcen, 6~étudiants actifs sur l'axe Tlemcen--Oran), répartis en deux groupes (9~passagers, 9~conducteurs).

\begin{table}[H]
\centering
\footnotesize
\begin{tabularx}{\textwidth}{|X|c|c|c|}
\hline
\rowcolor{rohlight}
\textbf{Groupe} & \textbf{N} & \textbf{Score SUS moyen} & \textbf{Interprétation} \\
\hline
Passagers (recherche \& réservation) & 9 & 74,2 / 100 & Bon (\textit{Good}) \\
\hline
Conducteurs (\textsc{kyc} \& publication) & 9 & 68,9 / 100 & Acceptable (\textit{OK}) \\
\hline
\rowcolor{rohlight!20}
\textbf{Score global} & \textbf{18} & \textbf{71,6 / 100} & \textbf{Bon (seuil~: 68)} \\
\hline
\end{tabularx}
  \caption{Résultats de l'évaluation SUS (18 participants, avril 2026)}
  \label{tab:sus}
\end{table}

L'échantillon, entièrement recruté en milieu universitaire, introduit un biais de sélection non négligeable qu'il convient d'analyser avec rigueur. Les 18~participants sont des étudiants en informatique, c'est-à-dire des \textit{digital natives} habitués à manipuler des interfaces mobiles complexes. Or, le segment cible de \textit{RohWinBghit} inclut également des conducteurs inter-wilayas âgés de 40 à 60~ans, dont le niveau de littératie numérique est structurellement inférieur. Le score SUS de 71,6 obtenu auprès d'un échantillon technophile doit donc être interprété comme une borne supérieure~: il est raisonnable d'anticiper un score significativement plus faible auprès de la population générale. Le score inférieur du groupe conducteurs (68,9 contre 74,2 pour les passagers) confirme cette hypothèse, la différence étant directement imputable à la complexité perçue du parcours \textsc{kyc} biométrique. Trois pistes concrètes d'amélioration \textsc{ux} ont été identifiées~: le traitement asynchrone du \textsc{kyc} en arrière-plan pour éliminer l'attente perçue, l'ajout de tutoriels animés (\textit{overlays}) pour guider les primo-utilisateurs, et l'adoption d'un profilage progressif permettant l'accès aux fonctions de base avant la complétion intégrale de la vérification. Une campagne d'évaluation élargie intégrant des profils démographiques variés est inscrite dans la feuille de route pré-commercialisation.

\textit{Retours qualitatifs~:} fluidité de la recherche, clarté du billet QR, lisibilité du dashboard conducteur. \textit{Axes d'amélioration~:} durée perçue du flux \textsc{kyc}, demande de mode sombre, retour d'état lors de connexion lente.

\subsection{Performances serveur}
\label{sec:performances-serveur}

Les mesures ont été conduites sur un VPS (4~vCPU, 8~Go~RAM) avec 50~utilisateurs concurrents simulés via k6.

\begin{table}[H]
\centering
\footnotesize
\begin{tabularx}{\textwidth}{|X|c|c|c|c|}
\hline
\rowcolor{rohlight}
\textbf{Route API} & \textbf{P50 (ms)} & \textbf{P95 (ms)} & \textbf{P99 (ms)} & \textbf{Taux succès} \\
\hline
\texttt{POST /auth/send-otp} & 48 & 112 & 198 & 99,8\,\% \\
\hline
\texttt{GET /trips/search} & 142 & 287 & 421 & 99,6\,\% \\
\hline
\texttt{POST /bookings} & 178 & 334 & 510 & 99,1\,\% \\
\hline
\texttt{GET /trips/:id/tracking} & 31 & 89 & 145 & 99,9\,\% \\
\hline
\texttt{POST /kyc/submit} & 1~240 & 2~800 & 4~100 & 98,2\,\% \\
\hline
WebSocket (message) & 12 & 38 & 72 & 99,9\,\% \\
\hline
\end{tabularx}
  \caption{Performances des routes API critiques (50 utilisateurs concurrents)}
  \label{tab:perf}
\end{table}

Les routes métier respectent l'objectif P95~$\leq$~500~ms, sauf le flux \textsc{kyc} (P95~=~2,8~s) lié au chargement du modèle InsightFace. Ce délai est communiqué via un indicateur de progression et une confirmation asynchrone.

\subsection{Algorithmes métier clés}

\paragraph{Matching conducteur--passager.} Le \texttt{MatchingService} évalue chaque trajet candidat par une fonction de score composite. Après application des critères éliminatoires (places disponibles, concordance de date, wilaya de départ et d'arrivée), les trajets restants sont classés selon le score pondéré~:
$$SC = 0{,}40 \times C_t + 0{,}35 \times Q_c + 0{,}15 \times R_{qp} + 0{,}10 \times B_{\text{kyc}}$$
où $C_t$ désigne la concordance temporelle (proximité entre l'heure de départ souhaitée et l'heure de départ réelle, normalisée sur $[0; 1]$), $Q_c$ la note moyenne du conducteur sur les 50~derniers trajets, $R_{qp}$ le ratio qualité/prix (inverse du prix par kilomètre, normalisé), et $B_{\text{kyc}}$ le bonus de vérification biométrique ($B_{\text{kyc}} = 1$ si statut \textit{verified}, $0$ sinon). Les coefficients ont été calibrés sur la base des résultats du questionnaire utilisateur (mars~2026), où la ponctualité et la confiance dans le conducteur constituaient les deux facteurs de décision prédominants.

\paragraph{Tarification dynamique (\textit{Surge Pricing}).} Un multiplicateur dynamique $m$ est calculé en temps réel pour chaque couple (wilaya, créneau horaire) selon la formule~:
$$m = \min\!\left(2{,}5\;; \; 1{,}0 + \alpha \cdot \frac{D_w - S_w}{\max(S_w, 1)}\right), \quad \alpha = 0{,}3$$
où $D_w$ représente la demande (nombre de recherches) et $S_w$ l'offre (nombre de places publiées) dans la fenêtre temporelle courante. Le coefficient d'amortissement $\alpha = 0{,}3$ garantit une montée progressive du multiplicateur, qui est borné supérieurement à $2{,}5$ pour préserver l'acceptabilité tarifaire~\cite{si2023what}. Le multiplicateur est mis en cache Redis avec un \textsc{ttl} de 5~minutes et affiché en transparence totale à l'utilisateur avant la confirmation de réservation.

\paragraph{Score conducteur.} Le score conducteur $SC_d \in [0; 100]$ est recalculé de manière asynchrone après chaque trajet via un \textit{worker} Bull/BullMQ, selon la décomposition~:
$$SC_d = 0{,}40 \times \bar{R} + 0{,}25 \times (1 - \tau_a) + 0{,}20 \times P + 0{,}10 \times A + 0{,}05 \times B_{\text{kyc}}$$
où $\bar{R}$ désigne la moyenne des notes passagers (normalisée sur $[0; 1]$), $\tau_a$ le taux d'annulation sur les 20~derniers trajets, $P$ le score de ponctualité (ratio des départs à l'heure), $A$ un bonus d'ancienneté (logarithmique en nombre de trajets), et $B_{\text{kyc}}$ le bonus de vérification biométrique.

% ================================================================
%  PARTIE II --- DÉPLOIEMENT
% ================================================================

\section{Déploiement}
\label{sec:deploiement}

Le backend Express.js est conteneurisé (Docker) et déployé sur un VPS Hetzner Cloud sous Ubuntu Server, derrière un reverse proxy Nginx assurant la terminaison TLS (Let's Encrypt) et l'équilibrage de charge. L'application Android est compilée via Expo Application Services (EAS) et disponible en test fermé sur Google Play Store. Le panneau d'administration React/Vite est servi par Nginx. Le monitoring est assuré par Sentry, les accès sécurisés par clés SSH et les journaux d'accès analysés en continu.

\subsection{Matrice de scalabilité de l'infrastructure}

Pour garantir la pérennité de la plateforme face à une croissance rapide, une stratégie de mise à l'échelle progressive a été définie (tableau~\ref{tab:scalability}).

\begin{table}[H]
\centering
\footnotesize
\begin{tabularx}{\textwidth}{|l|l|X|}
\hline
\rowcolor{rohlight}
\textbf{Phase (Utilisateurs)} & \textbf{Infrastructure} & \textbf{Topologie de la base de données} \\
\hline
\textbf{MVP ($<$ 5~000)} & VPS unique (8 vCPU, 16 Go RAM) & PostgreSQL local + sauvegardes distantes \\
\hline
\textbf{Croissance (5k -- 50k)} & Docker Swarm (3 nœuds) & PostgreSQL (Primary + 1 Read Replica) \\
\hline
\textbf{National ($>$ 500k)} & Cluster Kubernetes (K8s) dédié & Cluster PostgreSQL (Patroni/HAProxy) \\
\hline
\end{tabularx}
  \caption{Matrice d'évolution de l'infrastructure (Scalabilité)}
  \label{tab:scalability}
\end{table}

% ================================================================
%  PARTIE III --- ÉTUDE TECHNICO-ÉCONOMIQUE
% ================================================================

\section{Étude technico-économique}
\label{sec:business}

Ce mémoire s'inscrivant dans le cadre de l'\textbf{Arrêté interministériel n\textsuperscript{o}~1275} (mécanisme «~Un diplôme, une Startup~»), cette section présente l'étude de faisabilité économique.

\subsection{Conformité à l'Arrêté 1275}

\begin{table}[H]
\centering
\footnotesize
\renewcommand{\arraystretch}{1.3}
\begin{tabularx}{\textwidth}{|p{4cm}|X|p{2.5cm}|}
\hline
\rowcolor{rohlight}
\textbf{Critère} & \textbf{Justification} & \textbf{Statut} \\
\hline
Innovation technologique & Pipeline KYC biométrique, tarification dynamique, matching algorithmique & \textcolor{green!60!black}{\textbf{Conforme}} \\
\hline
Marché algérien identifié & 250~M voyages inter-wilayas/an, 1,9~M étudiants & \textcolor{green!60!black}{\textbf{Conforme}} \\
\hline
Équipe étudiante & 2 étudiants Master GL, Tlemcen & \textcolor{green!60!black}{\textbf{Conforme}} \\
\hline
Prototype fonctionnel & 62 écrans, 147/151 tests & \textcolor{green!60!black}{\textbf{Conforme}} \\
\hline
Étude économique & BMC, SWOT, prévisionnel 36 mois & \textcolor{green!60!black}{\textbf{Conforme}} \\
\hline
\end{tabularx}
  \caption{Conformité aux critères de l'Arrêté 1275}
  \label{tab:arrete1275}
\end{table}

\subsection{Analyse SWOT}

\begin{table}[H]
\centering
\begin{tabularx}{\textwidth}{|>{\columncolor{rohlight}}l|X|X|}
\hline
 & \textbf{Facteurs positifs} & \textbf{Facteurs négatifs} \\
\hline
\textbf{Interne} & \textbf{Forces~:} stack moderne, KYC biométrique intégré, couverture nationale, paiement local. & \textbf{Faiblesses~:} ressources financières limitées, absence de marque établie. \\
\hline
\textbf{Externe} & \textbf{Opportunités~:} segment inter-wilayas sous-servi, smartphone 75\,\%+, démographie jeune, adoption monétique +179\,\%. & \textbf{Menaces~:} extension de Yassir/Heetch, arrivée d'acteurs internationaux, dépendance à l'homologation SATIM. \\
\hline
\end{tabularx}
  \caption{Matrice SWOT du projet RohWinBghit}
\end{table}

\subsection{Modèle de monétisation}

Quatre sources de revenus sont prévues, par priorité de déploiement~:
\begin{enumerate}
  \item \textbf{Commission sur trajets (12\,\%)} prélevée sur chaque transaction confirmée.
  \item \textbf{Abonnement Premium conducteur (1~500~DZD/mois)}~: commissions réduites, badge Pro, carte thermique.
  \item \textbf{Publicité locale ciblée (in-app)} pour les annonceurs locaux.
  \item \textbf{Partenariats B2B} avec universités et entreprises.
\end{enumerate}

\subsection{Prévisionnel financier et Acquisition Client (CAC)}

Le modèle financier repose sur une ambition forte~: atteindre \textbf{50~000 utilisateurs actifs} à la fin de la première année. Cependant, cette croissance requiert une stratégie d'acquisition agressive. 

\paragraph{Le Paradoxe du CAC.} Avec un coût d'acquisition client (CAC) estimé à 350~DZD sur le marché numérique algérien, l'acquisition de 50~000 utilisateurs génère un besoin en financement marketing de \textbf{17~500~000~DZD}. Pour soutenir cette hyper-croissance et financer les coûts opérationnels de la première année (hébergement, salaires, frais juridiques), le projet nécessite impérativement une \textbf{levée de fonds (Seed) initiale de 20~000~000~DZD}.

\begin{table}[H]
\centering
\begin{tabularx}{\textwidth}{|X|r|r|r|}
\hline
\rowcolor{rohlight}
\textbf{Catégorie} & \textbf{Année~1 (DZD)} & \textbf{Année~2 (DZD)} & \textbf{Année~3 (DZD)} \\
\hline
\textbf{Revenus (Commissions + Premium)} & 6\,700\,000 & 29\,000\,000 & 92\,000\,000 \\
\hline
\textbf{Dépenses Opérationnelles (OPEX)} & 5\,200\,000 & 12\,500\,000 & 28\,000\,000 \\
\hline
\textbf{Dépenses Marketing (Acquisition)} & 17\,500\,000 & 8\,500\,000 & 14\,000\,000 \\
\hline
\rowcolor{rohlight!40}
\textbf{Résultat Net} & \textbf{-16\,000\,000} & \textbf{+8\,000\,000} & \textbf{+50\,000\,000} \\
\hline
\rowcolor{rohlight}
\textbf{Trésorerie (avec levée 20~M DZD)} & \textbf{+4\,000\,000} & \textbf{+12\,000\,000} & \textbf{+62\,000\,000} \\
\hline
\end{tabularx}
  \caption{Prévisionnel financier et flux de trésorerie sur 3 ans}
\end{table}

Ce modèle de capital-risque assume une perte nette la première année afin de capturer des parts de marché (stratégie de \textit{Blitzscaling}), avec un franchissement du \textbf{seuil de rentabilité (Break-Even) prévu au 2\textsuperscript{e} trimestre de l'Année~2}. La formule formelle d'évaluation des revenus de commission annuels est~:
$$R_n = U_n \times f \times p \times c$$
avec $U_n$ le nombre d'utilisateurs actifs, $f = 2{,}5$ la fréquence mensuelle d'utilisation (trajets/mois), $p = 800$~DZD le panier moyen par trajet, et $c = 0{,}12$ la commission de la plateforme. Pour l'Année~1 avec $U_1 = 50\,000$ utilisateurs actifs en régime permanent~:
$$R_1 = 50\,000 \times 2{,}5 \times 800 \times 0{,}12 \times 12 = 14\,400\,000 \text{~DZD (estimation haute)}$$
Le chiffre effectif de 6,7~M~DZD intègre la montée en charge progressive (l'objectif de 50~000 utilisateurs n'étant atteint qu'en fin d'année).

\subsection{Business Model Canvas}

\begin{table}[H]
\centering
\begin{tcolorbox}[
  enhanced,
  colback=white,
  colframe=rohblue,
  arc=6pt,
  boxrule=1.5pt,
  boxsep=2pt,
  left=2pt, right=2pt, top=4pt, bottom=4pt,
  toptitle=8pt, bottomtitle=8pt,
  title=\textbf{\large Business Model Canvas --- RohWinBghit},
  coltitle=white,
  colbacktitle=rohblue,
  center title
]
\renewcommand{\arraystretch}{1.6}
\noindent
\begin{tabularx}{\linewidth}{ >{\raggedright\arraybackslash}X | >{\raggedright\arraybackslash}X | >{\raggedright\arraybackslash}X | >{\raggedright\arraybackslash}X | >{\raggedright\arraybackslash}X }

\textcolor{rohblue}{\faHandshake\ \textbf{Partenaires}} &
\textcolor{rohblue}{\faCogs\ \textbf{Activités}} &
\textcolor{rohblue}{\faGem\ \textbf{Proposition de valeur}} &
\textcolor{rohblue}{\faHeart\ \textbf{Relation clients}} &
\textcolor{rohblue}{\faUsers\ \textbf{Segments}} \\[2pt]

\scriptsize
\begin{itemize}[leftmargin=0.8em, nosep]
  \item Mapbox, Firebase, Cloudinary
  \item Universités
  \item CIB/Edahabia (SATIM)
  \item Incubateurs (ANVREDET)
\end{itemize} &
\scriptsize
\begin{itemize}[leftmargin=0.8em, nosep]
  \item Développement mobile \& back-end
  \item Modération \& validation KYC
  \item Support client
\end{itemize} &
\scriptsize
\begin{itemize}[leftmargin=0.8em, nosep]
  \item Covoiturage inter-wilayas sécurisé
  \item KYC biométrique (IA)
  \item Économie partagée
  \item Paiement électronique local
\end{itemize} &
\scriptsize
\begin{itemize}[leftmargin=0.8em, nosep]
  \item Self-service (app intuitive)
  \item Confiance mutuelle (notation)
  \item Support réactif
\end{itemize} &
\scriptsize
\begin{itemize}[leftmargin=0.8em, nosep]
  \item Étudiants universitaires
  \item Travailleurs pendulaires
  \item Conducteurs (amortir coûts)
\end{itemize} \\[15pt]

\hline
\multicolumn{2}{p{\dimexpr0.4\linewidth-2\tabcolsep-1.2pt\relax}|}{\textcolor{rohblue}{\faCoins\ \textbf{Structure de coûts}}} &
\multicolumn{3}{p{\dimexpr0.6\linewidth-2\tabcolsep\relax}}{\textcolor{rohblue}{\faCreditCard\ \textbf{Sources de revenus}}} \\[2pt]

\multicolumn{2}{p{\dimexpr0.4\linewidth-2\tabcolsep-1.2pt\relax}|}{
\scriptsize
\begin{itemize}[leftmargin=0.8em, nosep]
  \item Masse salariale (tech/support)
  \item Hébergement VPS
  \item API payantes
  \item Acquisition client
\end{itemize}
} &
\multicolumn{3}{p{\dimexpr0.6\linewidth-2\tabcolsep\relax}}{
\scriptsize
\begin{itemize}[leftmargin=0.8em, nosep]
  \item Commission de 12\,\% sur trajets
  \item Abonnement Premium conducteurs
  \item Publicité in-app
  \item Partenariats B2B
\end{itemize}
} \\[6pt]

\end{tabularx}
\end{tcolorbox}
\caption{Business Model Canvas de RohWinBghit}
\end{table}

\subsection{Conformité juridique et Propriété Intellectuelle}

Sur les marchés numériques émergents, la conformité réglementaire absolue ne constitue pas seulement une obligation légale, mais un véritable \textbf{avantage concurrentiel (Moat)}. La plateforme \textit{RohWinBghit} est soumise à deux cadres majeurs~:

La \textbf{Loi~18-05 relative au commerce électronique}~\cite{loi1805_2018} impose la sécurisation stricte des paiements via le réseau interbancaire \textsc{satim}, l'établissement d'un contrat électronique par double clic, et la transparence tarifaire intégrale envers le consommateur. La \textbf{Loi~25-11 relative à la protection des personnes physiques dans le traitement des données à caractère personnel}~\cite{loi2511_2025} soumet la plateforme à l'autorité de l'\textbf{\textsc{anpdp}}. Le choix architectural fondamental d'isoler le moteur IA biométrique sur un serveur \textbf{souverain local} garantit qu'aucun vecteur facial de citoyen algérien ne transite par des \textsc{api} étrangères (Cloud Vision, AWS Rekognition). Couplé au chiffrement AES-256-GCM des plaques d'immatriculation au repos, cette souveraineté numérique bloque structurellement les acteurs internationaux non conformes souhaitant pénétrer le marché sans investissement en infrastructure locale.

\paragraph{Protection de la propriété intellectuelle.} Trois démarches de protection sont engagées parallèlement. Le dépôt de la marque \textit{RohWinBghit} auprès de l'\textbf{\textsc{inapi}} (Institut National Algérien de la Propriété Industrielle) confère un droit exclusif d'exploitation commerciale sur le territoire national pour une durée de dix ans renouvelable. Le scellement du code source auprès de l'\textbf{\textsc{onda}} (Office National des Droits d'Auteur et des Droits Voisins) établit une preuve d'antériorité horodatée, opposable en cas de litige de contrefaçon logicielle. Enfin, la combinaison de ces protections avec l'obtention du label startup dans le cadre de l'Arrêté~1275 constitue un triptyque juridique renforçant considérablement la crédibilité du projet auprès des investisseurs institutionnels lors de la levée de fonds Série~A envisagée en Phase~4 de la feuille de route.

\subsection{Feuille de route (18 mois)}
\label{sec:roadmap}

\begin{table}[H]
\centering
\footnotesize
\begin{tabularx}{\textwidth}{|l|l|X|}
\hline
\rowcolor{rohlight}
\textbf{Phase} & \textbf{Période} & \textbf{Objectifs clés} \\
\hline
\textbf{Phase 1} --- Fondation & Mois 1--3 & Finalisation MVP, dépôt CNL, création SARL, audit de sécurité. \\
\hline
\textbf{Phase 2} --- Bêta & Mois 4--9 & Déploiement axe Tlemcen--Oran--Alger, 1~000 utilisateurs, CIB/Edahabia en production. \\
\hline
\textbf{Phase 3} --- Expansion & Mois 10--15 & 20 wilayas, 10~000 utilisateurs, seuil de rentabilité, Premium. \\
\hline
\textbf{Phase 4} --- National & Mois 16--18 & 69 wilayas, 50~000 utilisateurs, préparation Série~A. \\
\hline
\end{tabularx}
  \caption{Feuille de route stratégique sur 18 mois}
  \label{tab:roadmap}
\end{table}

% ================================================================
%  PERSPECTIVES
% ================================================================

\section{Perspectives d'amélioration}
\label{sec:perspectives-amelioration}

Les perspectives d'évolution s'articulent autour de quatre axes~:
\begin{enumerate}
  \item \textbf{IA prédictive~:} intégration de modèles d'apprentissage automatique pour l'estimation de la demande et l'optimisation des coefficients de surge pricing.
  \item \textbf{Matching spatial avancé~:} algorithme sur graphe routier (PostGIS/OSRM) pour proposer des conducteurs transitant par la wilaya du passager, avec ajustement tarifaire du détour.
  \item \textbf{Intégration BaridiMob~:} pipeline OCR pour la validation automatique des reçus de transaction, élargissant l'inclusion financière.
  \item \textbf{Anti-fraude comportemental~:} détection des comptes multiples par empreinte matérielle et identification des schémas de réservation suspects par apprentissage non supervisé.
\end{enumerate}

\section{Conclusion}

Ce chapitre a exposé la validation de la plateforme (97,4\,\% de tests réussis, score SUS de 71,6/100, performances P95 conformes aux objectifs), son schéma de déploiement conteneurisé, et l'étude technico-économique confirmant la viabilité du modèle dans le cadre de l'Arrêté~1275.
